
Experiments expose a practical bottleneck in meta-learning for method selection: collecting comprehensive benchmark metadata from literature is harder than anticipated. Rather than refuting the meta-learning hypothesis, negative results identify an infrastructure gap that must be addressed before the approach can be tested fairly.

\textbf{Interpreting the Negative Results.} The core finding is that literature mining alone---querying APIs, parsing READMEs, extracting paper tables---provided only 14\% average coverage for critical dataset characteristics (sample\_size: 13.8\%, dimensionality: 0\%, class\_imbalance: 75.9\% but zero-variance artifact). This sparsity prevented testing whether dataset features correlate with method rankings (h-m1: zero computable correlations) and caused meta-classifier training to fail with degenerate input (h-m2: 1 usable feature, 25.6\% accuracy).

Why this is not hypothesis failure: all three sub-hypothesis deviations were classified as data limitation or scope change, not hypothesis issue. H-e1 successfully verified data source accessibility (OGB API loads, GitHub READMEs fetch, paper tables exist), achieving POC-level validation. H-m1 could not compute correlations because only 4 sample\_size values were extracted (insufficient for statistical testing), not because correlations were computed and found insignificant. H-m2 training failed because only 1 feature had sufficient coverage (num\_classes: 75.9\%), creating degenerate input that no meta-learning algorithm could learn from.

The real challenge is that meta-learning requires two-stage data collection. Stage 1 (identify benchmark sources) succeeded: OGB, FedML, LEAF, Champneys, and Zhou benchmarks are accessible via libraries, repositories, and publications. Stage 2 (extract dataset characteristics) is the bottleneck: APIs provide identifiers but not feature statistics, READMEs document usage but not sample sizes, and paper tables report method rankings but omit dataset properties. Extracting characteristics requires downloading raw datasets, loading them into memory, and computing statistics---an engineering effort beyond literature mining.

\textbf{Implications for Meta-Learning Research.} This work reveals that meta-learning research implicitly assumes dataset characteristics are available. AutoML systems (Auto-sklearn, TPOT) operate on datasets already loaded by practitioners, where extracting n, d, and class distribution is trivial. Literature-based meta-learning---training predictors from published benchmark results without dataset downloads---requires a metadata repository that does not currently exist.

Why literature mining falls short: published benchmarks prioritize documenting method comparisons (rankings, accuracy tables) over dataset properties. Zhou et al. report 9 medical imaging datasets with FedAvg/SCAFFOLD results but omit sample sizes and resolutions. Champneys et al. provide NARX-Poly versus LSTM comparisons but do not quantify input dimensionality or signal-to-noise ratios. This is rational for benchmark papers (readers care about ''which method wins'' not ''dataset n''), but it prevents aggregating metadata for meta-learning.

Proposed solution: a community-driven benchmark metadata repository analogous to Papers with Code leaderboards, but focused on dataset characteristics rather than method rankings. Such a repository would: (1) standardize feature extraction by defining Tier 1 universal features (sample\_size, dimensionality, class\_imbalance) and Tier 2 domain-specific features (autocorrelation, edge\_density, correlation\_rank), (2) automate collection by providing scripts to download benchmarks and compute features, reducing manual extraction errors (e.g., class\_imbalance artifact from template CSVs), and (3) link to method rankings by joining dataset characteristics with leaderboard results, enabling meta-learning training data generation. This infrastructure would unlock research questions currently untestable: Do dataset characteristics predict method performance? Can meta-classifiers generalize across domains? Are fast features sufficient, or do we need expensive probing?

\textbf{Honest Limitations.} Scope reduction: only 29 benchmarks were collected (versus 50-60 target), validating data source accessibility (POC level) but not executing exhaustive extraction. Full collection engineering (leaderboard scraping, Papers with Code authentication, PDF table parsing) was not pursued after discovering the metadata bottleneck. This scope reduction is principled: identifying the bottleneck early allowed analysis of root causes rather than investing effort in a flawed single-stage approach.

Manual extraction artifacts: manual CSV files for Champneys and Zhou benchmarks used standardized ranking percentiles [25, 50, 75, 100], producing zero-variance class\_imbalance (all values=0.559). This artifact reduced effective feature diversity. This is acceptable because manual extraction was a POC workaround for missing automated tools; it demonstrates the challenge (heterogeneous data formats) rather than invalidating findings. Real extraction from papers would require OCR or table parsing, confirming two-stage collection need.

Untested alternative explanations: the work did not test whether dataset downloads would yield complete metadata. It is possible that even with downloads, some benchmarks lack feature diversity (e.g., OGB ogbn-papers100M reports 111M samples, likely an API error). However, downloads would certainly improve coverage above 14\% (e.g., loading CIFAR-10 directly provides exact n=50K, d=3072, class distribution). Future work should quantify coverage improvement from two-stage collection.

Meta-learning hypothesis remains unverified: because h-m1 and h-m2 received insufficient data, the work cannot conclude whether dataset characteristics actually correlate with method rankings, or whether meta-classifiers can learn these relationships. The hypothesis is untested, not disproven. Negative results should be interpreted as ''literature mining insufficient'' rather than ''meta-learning does not work.''

\textbf{Methodological Contributions.} Despite negative results on the core hypothesis, this work contributes three methodological insights. First, transparent negative result reporting documents failure modes with quantitative evidence (14\% average feature coverage, zero computable correlations, 1 usable feature after preprocessing) and classifies deviations by type (data limitation versus hypothesis issue). This transparency allows readers to assess whether failures stem from experimental execution or theoretical invalidity. Mock data elimination protocols (external review, variance checks, coverage reporting) verify real data usage, preventing false negatives from hard-coded defaults.

Second, staged hypothesis testing by decomposing the core hypothesis into h-e1 (data collection) $\rightarrow$ h-m1 (correlation) $\rightarrow$ h-m2 (meta-learning) isolates failure modes. When h-e1 achieves POC validation but h-m1 fails, the issue is feature extraction (not source availability). When h-m1 shows zero correlations due to insufficient data (not insignificant p-values), the issue is input quality (not absence of relationships). Staged testing prevents misattribution: distinguishing ''not enough data'' from ''data exists but is uninformative'' from ''features are informative but classifier does not learn.''

Third, infrastructure gap identification frames negative results as a practical bottleneck (metadata extraction) rather than hypothesis refutation (meta-learning does not work), which is actionable. Findings suggest where research effort should focus: building automated extraction tools and metadata repositories, not refining meta-learning algorithms. This reframes the research question from ''does meta-learning predict method selection?'' to ''what infrastructure enables fair testing of meta-learning?''

\textbf{Broader Impact.} For practitioners, do not expect literature mining (reading papers, checking leaderboards) to provide sufficient dataset characteristics for informed method selection. If dataset properties are needed to guide method choice, download and analyze raw data rather than relying on published metadata.

For researchers, when proposing meta-learning approaches, verify that training data (dataset characteristics) is accessible. If an approach requires feature X, check whether X is reported in literature or must be computed from raw datasets. This work shows that seemingly ''trivial'' features like sample\_size and dimensionality are absent from 86-100\% of published benchmarks.

For benchmark publishers, consider documenting dataset characteristics alongside method rankings. Including a ''Dataset Properties'' section (n, d, class distribution, domain-specific statistics) in benchmark papers would enable meta-learning research without requiring dataset downloads. Standardizing reporting (e.g., mandatory fields in benchmark submission templates) could accelerate meta-learning adoption.

\textbf{Future Work.} Immediate next step: implement Stage 2 data collection for the 29 collected benchmarks. Download OGB datasets, FedML repositories, and benchmark datasets from Champneys/Zhou sources. Compute Tier 1 features from raw data (not API metadata). Re-run h-m1 correlation analysis with complete features (expected: at least 10 benchmarks with full Tier 1 coverage) to test whether correlations emerge with richer data.

Medium-term goal: extend collection to 50-60 benchmarks with automated extraction tools (leaderboard scrapers, OCR for paper tables, dataset loaders for OGB/FedML/LEAF). Test h-m2 meta-classifier training with at least 40 complete feature vectors and at least 5 informative features. Compare two-stage (downloads + extraction) versus single-stage (literature mining only) coverage and meta-learning performance.

Long-term vision: build a benchmark metadata repository integrating dataset characteristics with Papers with Code leaderboards. Community-contributed extraction scripts ensure standardization. Researchers train meta-classifiers on aggregated data without per-project manual collection. Enable new research: temporal generalization (do feature-method correlations hold across eras?), domain transfer (do patterns from vision generalize to NLP?), minimal feature sets (which Tier 1 features are sufficient?).

Alternative approaches: if even two-stage collection proves insufficient (e.g., some benchmarks lack raw data), explore Tier 3 features (model probing via quick training runs on new datasets). Tier 3 is expensive (approximately 5-10 GPU-minutes per dataset) but guarantees feature availability. Trade-off: violates ''fast computation'' constraint (under 1 minute) but may be necessary for fair meta-learning testing.
